%% ==========================================================================
%%  5  Solved Cases
%%
%%  Theorem C (binary additive) and Theorem D (identical costs) of the working
%%  note.  Both proofs work by making the arc-weight bound telescope along
%%  paths; that mechanism is the thing to generalise.
%% ==========================================================================
\section{Solved Cases}
\label{sec:solved}

Conjecture~\ref{conj:main} holds in three restrictions of the model, by explicit
polynomial-time constructions, and on one further class of instances cut out by
a certificate rather than by a restriction on the valuations
(\Cref{sec:balancedclass}). The first three proofs run the same way: bound
$\edgew{\alloc}(i,j)$ by a difference $\phi_i - \phi_j$ of a per-agent potential,
so that the bound telescopes to $\phi_{i_1} - \phi_{i_r}$ along a path and to
$0$ around a cycle. Envy-freeability and the $\set{0,1}$ bound then both follow
from controlling the spread of $\phi$. Reproducing that mechanism in general is
the open problem.

\paragraph{What is and is not new here.}
This section should be read with its overlap against the literature stated
plainly, because two of the three cases are largely or wholly not ours.

\begin{itemize}
  \item \textbf{Binary additive costs} (Theorem~\ref{thm:binadd}) are
    \emph{subsumed} by~\cite{LMS26}, which settles all additive goods-and-chores
    instances. See Remark~\ref{rem:binadd-superseded}. The statement is not new;
    only the potential-and-telescoping mechanism is retained, because that is
    what the general case has to reproduce.
  \item \textbf{Identical costs} (Theorem~\ref{thm:identical}) are \emph{not}
    subsumed, for the reason given in Remark~\ref{rem:identical-new} --- but the
    result is elementary and is best read as a warm-up.
  \item \textbf{Two agents} is~\cite[Theorem~4]{BKNS22}. Only the reduction
    carrying it to chores (Remark~\ref{rem:complement}) is ours.
  \item \textbf{The uniformly balanced class} (\Cref{sec:balancedclass}) is new,
    and is the only case here not implied by a cited paper. It is a certificate
    rather than an algorithm, and it is incomplete.
\end{itemize}

\noindent
Nothing in this section, in other words, is load-bearing. The substance of the
paper is in \Cref{sec:approach1,sec:approach2,sec:approach4,sec:approach5},
which are negative, and in \Cref{sec:structure,sec:approach3}, which are
structural.

% --------------------------------------------------------------------------
\subsection{Binary additive costs}
\label{sec:binadd}

Let $\cost_i(S) = \abs{S \cap D_i}$ for a set $D_i \subseteq \items$ of chores
that agent $i$ dislikes; every other chore is free for her. Call a chore
\emph{universally bad} if it lies in $D_i$ for every $i$, write $U$ for the set
of universally bad chores and $q := \abs{U}$.

\begin{theorem}
\label{thm:binadd}
For binary additive costs, the following allocation is envy-freeable with
minimum subsidy $\optsubsidy \in \set{0,1}^n$, hence total subsidy at most
$n-1$, and it is computable in time $O(nm)$.
\begin{quote}
Give each chore $g \notin U$ to some agent $i$ with $g \notin D_i$; at least one
exists, by definition of $U$. Distribute $U$ as evenly as possible, so that
$u_i := \abs{\bundle{i} \cap U} \in \set{\lfloor q/n \rfloor, \lceil q/n \rceil}$
for every $i$.
\end{quote}
The bound is tight, by Example~\ref{ex:lowerbound}.
\end{theorem}

\begin{proof}
By construction every chore in $\bundle{i} \setminus U$ is free for $i$, so
$\cost_i(\bundle{i}) = \abs{\bundle{i} \cap D_i} = u_i$. For any $j$ we have
$\bundle{j} \cap U \subseteq \bundle{j} \cap D_i$, since $U \subseteq D_i$, and
therefore $\cost_i(\bundle{j}) = \abs{\bundle{j} \cap D_i} \ge u_j$. Hence
\[
  \edgew{\alloc}(i,j) \;=\; \cost_i(\bundle{i}) - \cost_i(\bundle{j})
  \;\le\; u_i - u_j .
\]
Summing along a path $P = (i_1, \dots, i_r)$ the right-hand side telescopes:
\[
  \edgew{\alloc}(P) \;\le\; u_{i_1} - u_{i_r}
  \;\le\; \lceil q/n \rceil - \lfloor q/n \rfloor \;\le\; 1 .
\]
Around a cycle the same computation gives $\edgew{\alloc}(C) \le 0$, so
$\envygraph{\alloc}$ has no positive-weight cycle and $\alloc$ is envy-freeable
by Theorem~\ref{thm:hs-characterisation}. By Theorem~\ref{thm:hs-minsubsidy},
$\optsubsidy_i = \pathw{\alloc}(i) \le 1$, and
Observation~\ref{obs:integrality} forces $\optsubsidy \in \set{0,1}^n$. Since
some agent receives $\lfloor q/n \rfloor$ universally bad chores and hence
requires no subsidy, the total is at most $n-1$. Example~\ref{ex:lowerbound} is
binary additive with $D_i = \items$ for every $i$, so the bound cannot be
improved.
\end{proof}

\begin{remark}[Superseded, and by how much]
\label{rem:binadd-superseded}
Theorem~\ref{thm:binadd} is subsumed by~\cite{LMS26}, which appeared in July
2026 and gives one dollar per agent for \emph{all} additive goods-and-chores
instances with $\abs{v_i(g)} \le 1$; binary additive costs are the special case
$v_i(g) \in \set{-1,0}$. We keep the proof because the mechanism, not the
statement, is what the general case needs: the potential $\phi_i = u_i$ and its
telescoping along paths is the only device in this paper that produces a
$\set{0,1}$ bound from scratch, and the general conjecture is precisely the
problem of finding the right $\phi$. Theorem~\ref{thm:identical} and the $n = 2$
case are \emph{not} subsumed, since dichotomous costs need not be additive.
\end{remark}

The construction is a two-phase version of the shape the general case must
reproduce: \emph{dump every chore on somebody who finds it free, then balance
whatever nobody wants}. The obstruction to running the same argument for general
negative dichotomous costs is that ``free for $i$'' becomes context dependent. A
chore may be free on top of $\bundle{j}$ and cost a full unit on top of
$\bundle{i}$, so there is no fixed set $D_i$ to compute with and no fixed
potential $u_i$ to telescope against. Theorem~\ref{thm:obstruction} shows that
this is a genuine obstacle rather than a bookkeeping nuisance.

% --------------------------------------------------------------------------
\subsection{Identical costs, and two agents}
\label{sec:identical}

\begin{theorem}
\label{thm:identical}
Suppose $\cost_1 = \dots = \cost_n = \cost$. Then every allocation is
envy-freeable, and greedily adding each chore to a currently cheapest bundle
yields $\optsubsidy \in \set{0,1}^n$ in time $O(m(n + \tau))$, where $\tau$ is
the cost of one oracle call.
\end{theorem}

\begin{proof}
For any permutation $\sigma$ we have $\sum_i \cost(A_{\sigma(i)}) =
\sum_i \cost(\bundle{i})$, since both sides sum the same function over the same
multiset of bundles. So condition~(ii) of Theorem~\ref{thm:hs-characterisation}
holds with equality for every allocation, and every allocation is
envy-freeable. Moreover $\edgew{\alloc}(i,j) = \cost(\bundle{i}) -
\cost(\bundle{j})$ telescopes \emph{exactly} along a path
$P = (i_1,\dots,i_r)$, giving $\edgew{\alloc}(P) = \cost(A_{i_1}) -
\cost(A_{i_r})$ and hence
\[
  \pathw{\alloc}(u) \;=\; \cost(\bundle{u}) - \min_{j \in \agents}
  \cost(\bundle{j}) .
\]
It therefore suffices to keep $\max_i \cost(\bundle{i}) - \min_i
\cost(\bundle{i}) \le 1$. This is an invariant of the greedy rule. It holds
vacuously at the empty allocation. Suppose it holds before a step, so that
$\cost(\bundle{i}) \in \set{\mu, \mu+1}$ for every $i$, where $\mu$ is the
current minimum, and let the next chore go to a bundle attaining $\mu$. Its cost
becomes $\mu$ or $\mu + 1$, because marginals lie in $\set{0,1}$; every other
bundle is unchanged. So all costs still lie in $\set{\mu,\mu+1}$ and the
invariant survives. The bound $\optsubsidy \in \set{0,1}^n$ follows with
Observation~\ref{obs:integrality}.
\end{proof}

\begin{remark}[Why this one survives the literature]
\label{rem:identical-new}
Theorem~\ref{thm:identical} is not covered by~\cite{LMS26}, because identical
dichotomous costs need not be additive. The instance $\cost_i(S) =
\max(0,\abs{S}-1)$ for every $i$ is identical and dichotomous, with marginals
$0,1,1,\dots$, and is supermodular rather than additive --- it is the very cost
function that drives the obstruction of Theorem~\ref{thm:obstruction}. Nor
does~\cite{BKNS22} cover it, that being a goods result, and~\cite{KMSTY24} gives
only $n-1$ per agent on this class. So the statement is new, in the narrow sense
that no cited paper implies it.

It is also easy, and should not be oversold. With identical costs every agent
agrees on the cost of every bundle, the problem degenerates to load balancing,
and the potential $\phi_i = \cost(\bundle{i})$ telescopes \emph{exactly} rather
than merely bounding the arcs. That exactness is what makes the proof three
lines long, and it is precisely what fails as soon as two agents disagree.
\end{remark}

For $n = 2$ the conjecture holds as well, with total subsidy at most $1$, but
here the theorem is somebody else's: by Remark~\ref{rem:complement} two-agent
chore division \emph{is} two-agent goods division with the bundles swapped, so
the content is~\cite[Theorem~4]{BKNS22} and ours is only the reduction. As noted
there, that reduction consumes the identity $\items \setminus \bundle{j} =
\bundle{i}$ and so says nothing for $n \ge 3$.

% --------------------------------------------------------------------------
\subsection{Small bundles, and fewer chores than agents}
\label{sec:smallbundle}

Every other case in this section, and every approach in the paper, eventually
needs an \emph{exchange} step: given a bad allocation, produce a better one by
moving items. That is exactly what dichotomous costs refuse to supply, since a
bundle of positive cost need not contain any single item whose removal lowers it
--- for instance $\cost(S) = \min(\max(0,\abs{S}-1),1)$ on a three-element set
has $\cost = 1$ while every single removal leaves it at $1$. The following
theorem needs no exchange at all; it reads the conclusion straight off the
cycle-closing bound.

\begin{theorem}[Small-bundle theorem]
\label{thm:smallbundle}
Let $\alloc$ be envy-freeable and suppose every bundle costs at most one unit to
every agent: $\cost_v(\bundle{i}) \le 1$ for all $v, i \in \agents$. Then
$\optsubsidy \in \set{0,1}^n$ and the total subsidy is at most $n-1$.
\end{theorem}

\begin{proof}
By Theorem~\ref{thm:cyclebound},
$\pathw{\alloc}(u) \le \max\big(0, \max_{v \ne u}[\cost_v(\bundle{u}) -
\cost_v(\bundle{v})]\big) \le \max(0, 1 - 0) = 1$, using
$\cost_v(\bundle{u}) \le 1$ and $\cost_v(\bundle{v}) \ge 0$.
Observation~\ref{obs:integrality} then gives
$\optsubsidy \in \set{0,1}^n$, and the endpoint of a heaviest path has
$\pathw{} = 0$, so the total is at most $n-1$.
\end{proof}

\begin{corollary}[Fewer chores than agents]
\label{cor:m-le-n}
If $m \le n$ then Conjecture~\ref{conj:main} holds, and a witness is computable
by a single minimum-cost bipartite matching.
\end{corollary}

\begin{proof}
Among all allocations giving each agent at most one chore, take one minimising
$\sum_i \cost_i(\bundle{i})$; such an allocation exists because $m \le n$, and
it is a minimum-cost matching between agents and chores. Every permutation of
its own bundles is again an allocation of this form, so the minimisation was
over a set containing all of them, and condition~(ii) of
Theorem~\ref{thm:hs-characterisation} holds: $\alloc$ is envy-freeable. Each
bundle has $\abs{\bundle{i}} \le 1$, so $\cost_v(\bundle{i}) \le
\abs{\bundle{i}} \le 1$ for every $v$ --- because marginals lie in $\set{0,1}$
--- and Theorem~\ref{thm:smallbundle} applies. The bound is tight by
Example~\ref{ex:lowerbound}, which has $m = n-1$.
\end{proof}

The hypothesis of Theorem~\ref{thm:smallbundle} is genuinely weaker than
$m \le n$: it also covers instances with many chores in which every agent finds
all but a few of them free, so that no bundle is expensive to anybody. Both
statements were checked semantically against the true longest-path computation
--- exhaustively for $m \le 3$, $n \le 3$, on $3{,}000$ random instances for the
theorem, and on $111{,}232$ instances for the corollary --- with no violations
(\texttt{structural.py}).

\begin{remark}[Why this is worth stating despite being easy]
\label{rem:smallbundle-scope}
The proof is three lines and the class is small, so the content is not the
theorem but the \emph{mechanism}: it is the only argument in this paper that
reaches $\pathw{} \le 1$ without an exchange step, a balance hypothesis, a
schedule, or a search. It does so by making the cycle-closing bound tight from
the far side --- bounding $\cost_v(\bundle{u})$ from above rather than
$\cost_v(\bundle{v})$ from below. Whether that trick has a wider reach is open;
the obvious next question is what happens when bundles cost at most $2$, where
the same computation gives only $\pathw{} \le 2$ and the argument stops.
\end{remark}

% --------------------------------------------------------------------------
\subsection{Instances admitting a uniformly balanced family}
\label{sec:balancedclass}

The fourth case is not a restriction on the valuation class but on the instance,
and it is the only one of the four that is genuinely new here.

\begin{theorem}
\label{thm:balanced-class}
Conjecture~\ref{conj:main} holds for every instance admitting a uniformly
balanced family, i.e.\ a partition of $\items$ into $n$ bundles that every agent
values within one unit of each other
(\Cref{def:unifbal}). Given such a family, an envy-free solution with
$\optsubsidy \in \set{0,1}^n$ and total subsidy at most $n-1$ is obtained by
assigning the bundles by a single maximum-weight matching.
\end{theorem}

\begin{proof}
Proposition~\ref{prop:balance-suffices}.
\end{proof}

Three things should be said about the scope of this, none of which the proof
makes evident.

\begin{itemize}
  \item \textbf{It is a certificate, not an algorithm.} Verifying a proposed
    family and assigning it are polynomial; \emph{finding} one is a search over
    partitions, and no polynomial procedure for that is known. So
    Theorem~\ref{thm:balanced-class} decides an instance only once somebody
    supplies the family.
  \item \textbf{The class is large but not everything.} It contains every
    instance at $n = m = 3$ --- all $9{,}880$, checked exhaustively --- and every
    hard instance the project had accumulated before it was tested. It does
    \emph{not} contain the discrepancy instance of
    Proposition~\ref{prop:no-balance}, at $n = 3$ and $m = 4$, nor the
    certified-hard set-splitting instance.
  \item \textbf{Membership is not the same as difficulty.} The instances outside
    the class are not the instances that need large subsidies; on the
    counterexample of Proposition~\ref{prop:no-balance}, $19$ allocations are
    good and several need no subsidy at all. Uniform balance is a property of
    the certificate, not of the instance's hardness.
\end{itemize}

So this is a solved case in the same sense as the others in this section: real,
checkable, and not load-bearing. What it contributes is the counterexample and
Remark~\ref{rem:arc-vs-path}, which say that any complete method must reason
about paths rather than arcs.

% --------------------------------------------------------------------------
\subsection{The residual at \texorpdfstring{$n = 3$}{n=3}, and who gets paid}
\label{sec:n3-residual}

Every approach in this report has attacked the general case, and each died on a
general obstruction. The one gap in that coverage is a dedicated attack on a
bounded case beyond $n = 2$. At $n = 3$ the target is finite in a useful way: by
Proposition~\ref{prop:zero-coordinate} Conjecture~\ref{conj:main} is the single
clause $\max_i \pathw{\alloc}(i) \le 1$, and by the two-tier characterisation with
Theorem~\ref{thm:paidsets-lattice} the paid set lies in a lattice inside
$2^{\set{1,2,3}}$ in which $\emptyset$ and $\agents$ impose the same condition. So
there are exactly three cases, $\abs{S} = 0, 1, 2$, and
Example~\ref{ex:lowerbound} shows the last is sometimes forced.

\begin{observation}[The residual is small and its paid set is extremal]
\label{obs:n3-residual}
Over $2{,}460$ instances at $n = 3$ drawn from six families --- uniform, nested,
mixed, capped, threshold and disjoint-interest --- searched exhaustively over all
$3^m$ allocations:
\begin{itemize}
  \item every instance admits a good allocation; the minimal paid-set size is
    $0$ on $2{,}304$, $1$ on $135$ and $2$ on $21$, so the residual beyond exact
    envy-freeness is $156$ instances, $6.2\%$;
  \item in \emph{every} minimal witness with $\abs{S} = 1$, the paid agent attains
    $\max_i \cost_i(\bundle{i})$ --- $135$ of $135$;
  \item in \emph{every} minimal witness with $\abs{S} = 2$, the unpaid agent
    attains $\min_i \cost_i(\bundle{i})$ --- $21$ of $21$.
\end{itemize}
Bundle size does \emph{not} characterise it: the paid agent holds a largest
bundle in only $91$ of $135$ witnesses, and the unpaid agent a smallest in $2$ of
$21$ (\texttt{update\_44/n3\_cases.py}).
\end{observation}

So at a minimal witness the singled-out agent sits at an extreme of \emph{own
cost}, which removes the choice of \emph{which} agent is paid: one may assume it
is an $\argmax_i \cost_i(\bundle{i})$ when $\abs{S} = 1$, and that the unpaid one
is an $\argmin$ when $\abs{S} = 2$. That is a necessary condition, not a
sufficient one, and it is what a case analysis at $n = 3$ would start from.

\begin{remark}[The tempting unification is false]
\label{rem:n3-unified-false}
Both clauses look like instances of $S = \set{i : \cost_i(\bundle{i}) >
\min_j \cost_j(\bundle{j})}$ --- the paid agents being exactly those not attaining
the minimum own cost. They are not. Over the residual instances that equality
holds at \emph{some} minimal witness on $59$ of $69$ at $n = 3$, and at
\emph{every} one on only $37$; the failure is visible already when two agents tie
at the maximum while $\abs{S} = 1$. At $n = 4$ and $n = 5$ it fares better
($48/53$ and $67/69$ for every witness), which is a warning rather than a
comfort: the equality is not a theorem and should not be assumed.
\end{remark}

\begin{remark}[No counterexample, and the limits of looking]
\label{rem:no-counterexample}
Conjecture~\ref{conj:main} was also attacked directly, by searching for an
instance admitting \emph{no} good allocation. Over $683$ instances at
$n \le 6$, $m \le 11$, each searched exhaustively over all $n^m$ allocations and
drawn from the adversarial families above, none was found. The power of that
search is limited and should be read as such: only $32$ of the $683$ were hard at
all, in the sense of admitting no exactly envy-free allocation, so the conjecture
was under strain on under $5\%$ of the sample
(\texttt{update\_44/counterexample\_hunt.py}).
\end{remark}

\begin{remark}[Why the two-agent theorem does not bootstrap]
\label{rem:n2-no-bootstrap}
In the $\abs{S} = 1$ case the two unpaid agents are exactly envy-free between
themselves, which invites using Theorem~\ref{thm:n2-complete} on the sub-instance
they share. It does not apply: that theorem delivers a subsidy of at most one for
two agents, not exact envy-freeness, and Example~\ref{ex:lowerbound} at $n = 2$
already forces a dollar. The sub-problem at $n = 3$ is strictly harder than the
one solved at $n = 2$, and the case must be attacked directly.
\end{remark}
